\subsection{USE OF COUNTABLE SEQUENCES}

Proposition 6 is the origin of the part played by countable sequences of points in the theory of metrizable spaces; for many problems, they can be used to advantage in place of filters. This is because the neighbourhood filters of points of a metrizable space (and therefore also convergent filters) are determined by convergent sequences of points of the space: for since the neighbourhood filter of a point has a countable base, it is the intersection of the elementary filters finer than itself (Chapter I, § 6, no. 8, Proposition 11), i.e., of the elementary filters associated with sequences which converge to the point in question.

On the other hand, the notion of a convergent sequence is not adapted to the study of topological spaces in which there are points whose neighbourhood filter has no countable base. In particular, Hausdorff non-discrete topological spaces can be constructed in which, at every point x, the intersection of any countable family of neighbourhoods of x is again a neighbourhood of \(x (*)\) ; in such a space the only convergent sequences are those in which all the terms are equal from a certain index onwards.

As examples of the use of countable sequences we give the following propositions:

PROPOSITION 8. In a metrizable space X, a point x lies in the closure of a non-empty subset A of X if and only if there is a sequence of points of A which converges to x.

We know already, from Chapter I, § 7, no. 3, that the condition is sufficient. To see that it is necessary, consider a countable fundamental system \((\mathbf{V}_n)\) of neighbourhoods of \(x\) such that \(\mathbf{V}_{n+1} \subset \mathbf{V}_n\) for each \(n\) . If \(x\) lies in the closure of A then each \(\mathbf{V}_n\) meets A, and if \(x_n\) lies in \(\mathbf{V}_n \cap \mathbf{A}\) , the sequence \((x_n)\) converges to \(x\) (**).

From Proposition 8 we deduce:

PROPOSITION 9. A metric space X is complete if and only if every Cauchy sequence in X is convergent.

Let \(\hat{\mathbf{X}}\) be the completion of \(\mathbf{X}\) . If there is a point \(x \in \hat{\mathbf{X}}\) which does not belong to \(\mathbf{X}\) , then there is a sequence \((x_n)\) of points of \(\mathbf{X}\) which converges to \(x\) , and this is a non-convergent Cauchy sequence in \(\mathbf{X}\) .

PROPOSITION 10. Let X be a metrizable space and let f be a mapping of X into a topological space \(X'\) . Then f is continuous at a point \(x \in X\) if and only if, whenever \((x_n)\) is a sequence of points of X which converges to x, the sequence \((f(x_n))\) converges to \(f(x)\) in \(X'\) .

The condition is necessary, from Chapter I, § 7, no. 4, Proposition 9, Corollary 1. To show that it is sufficient, consider the filter \(\mathfrak{B}'\) of neighbourhoods of \(f(a)\) in \(\mathbf{X}'\) ; the hypothesis implies that \(\overline{f}(\mathfrak{B}')\) is coarser than every elementary filter associated with a sequence which converges to \(a\) , that is to say every elementary filter which converges to \(a\) ; but the intersection of these filters is the neighbourhood filter of \(a\) (Chapter I, § 6, no. 8, Proposition 11). Hence the result.

Note that Propositions 8 and 10 are valid in any space X in which every point has a countable fundamental system of neighbourhoods.

